\section{Materials and methods}\label{sec:methods}

\subsection{Study design and analysis plans}

This is a retrospective modelling study on public, de-identified fundus photographs. An
evaluation protocol (unit of independence, held-out sets, model selection on development data
only, metrics and uncertainty) was fixed before the first training run. For the analyses
reported here, the plan --- hypothesis, control, analysis and the result that would falsify
it --- was committed to the public repository before the analysis was run: the derivation test
of the macular-oedema grade (Supplementary Note~2), the nine journal outcomes O1--O9 (seed
replicates, reverse direction, second external corpus, five-fold foundation model) and the
revision outcomes A1--E2 (mechanism, label-free rules, sample size, two further external sets
and the models trained on EyePACS). These are analysis plans prospectively committed to version
control, not entries in a public registry; their commit identifiers are listed in Supplementary
Table~S1. The repository history begins on 31 August 2026, so the pre-specification of
development runs trained before that date cannot be verified from it, and we do not rely on it.
Reporting follows TRIPOD+AI\cite{collins2024tripodai} and CLAIM\cite{tejani2024claim}
(Supplementary Table~S2).

\subsection{Data}

The \emph{development pool} combined IDRiD\cite{porwal2018idrid,porwal2020challenge}
(\pn{n.idrid} images) and Messidor-2\cite{decenciere2014messidor,abramoff2013messidor2}
(\pn{n.messidor2} gradable images with adjudicated grades\cite{krause2018grader}),
\pn{n.dev} images in total (Table~\ref{tab:cohorts}; Fig.~\ref{fig:design}). The external sets, never used for
training, model selection or threshold fitting, were APTOS 2019 (\pn{n.aptos} images, India,
single grader), DDR\cite{li2019ddr} (\pn{n.ddr} gradable images, 147 hospitals in China,
seven graders, distributed at 512 pixels), the regular fundus images of
DeepDRiD\cite{liu2022deepdrid} (\pn{r2.deepdrid.n} images of 500 patients, Shanghai) and
SUSTech-SYSU\cite{lin2020sustech} (\pn{r2.sustech.n} images with an international-scale grade,
Shenzhen). In the reverse direction APTOS was the training corpus and Messidor-2 and IDRiD were
external. The EyePACS 2015 training set (\pn{n.eyepacs} images, single grader) was used to
pretrain the development models and, at revision, as the development set of two further models.
All corpora use the international 5-grade scale\cite{wilkinson2003scale}; referable DR is grade
$\geq$2 and, because macular oedema is not graded in the external sets, it is defined by DR
grade alone. A content audit (256-bit difference hash after the training crop) found no
development image within re-encoding distance of any external image (Supplementary Note~7).
Patient identifiers are published only for DeepDRiD and EyePACS; Messidor-2 is distributed as
pairs of images per examination, but the pairing file is not part of the mirrors used, so
resampling elsewhere is by image (near-duplicates merged).

\subsection{Models}

All models share one design: a backbone, global average pooling and ordinal threshold heads
that output four logits for $P(y>k)$\cite{cao2020rank}, trained with masked binary
cross-entropy, a grade-balanced sampler, AdamW, cosine decay, weight averaging and
dihedral test-time augmentation (details in Supplementary Table~S3). The reference model is a
DenseNet121\cite{huang2017densenet} pretrained on ImageNet and then EyePACS and fine-tuned at
448 pixels; the selected model replaces the backbone with EfficientNet-B3\cite{tan2019efficientnet}.
The other eight pre-registered development models change one factor each
(Table~\ref{tab:models}): resolution (224, 560 and native 640 pixels), a macula-centred crop,
linear-probe-then-fine-tune\cite{kumar2022lpft}, CORAL and CORN heads\cite{cao2020rank,shi2023corn}
and ConvNeXt-tiny\cite{liu2022convnext}; the retinal foundation model RETFound\cite{zhou2023retfound}
was fine-tuned separately. Development results are out of fold over a frozen 5-fold split
stratified by corpus and grade; external predictions average the five fold models. The models
trained on EyePACS (DenseNet121 and EfficientNet-B3, 448 pixels, 15 epochs from ImageNet) used a
patient-grouped split, trained on four folds and used the fifth (\pn{eye.dn.devn} images) for
checkpoint selection and for fitting development rules.

\subsection{Decision rules and the matched decision rule}

For an ordinal model with $p_k=P(y>k)$ the expected grade is $s=\sum_k p_k$. Every model's
outputs were decoded under the same rules (Supplementary Note~1). \textbf{R1} is the shipped
rule: the number of leading thresholds with $p_k>0.5$. \textbf{R2} places four cut-points on $s$
to maximise QWK on development out-of-fold predictions, and \textbf{R3} moves R2's grade-1/2
cut to give 90\% development referable sensitivity. \textbf{R4} and \textbf{R4m} fit the
cut-points on the external labels, maximising QWK or macro-averaged recall, by two-fold
cross-fitting; they are diagnostics of what local labels could recover, never performance.
\textbf{R5} adapts the class posteriors (derived from the ordinal outputs) to the external
population without labels by EM estimation of the target prior\cite{saerens2002adjusting},
and \textbf{R5b} by black-box shift estimation\cite{lipton2018bbse}; both decode by argmax.

The \emph{matched decision rule} is the comparison rule used throughout: no difference between
two models is attributed to the representation unless it survives refitting both models'
cut-points by the same cross-fitted procedure (on development data, fold $f$'s cut-points come
from the other four folds; externally, R4). Because a per-grade verdict can depend on how the
cut-points are cross-fitted (Supplementary Note~3), every matched claim names its procedure.

\subsection{Decomposing external performance}

Discrimination was measured at every ordinal threshold (area under the ROC curve of $p_k$ for
$y>k$, $k=0,\dots,3$) and overall by Harrell's C of $s$. The mapping from score to grade was
described by the mean expected grade given the true grade, $\mathrm{E}[s\mid y=k]$; under a pure
prior (label) shift this quantity is the same in every population, so a grade-by-grade departure
from its development value indicates a conditional shift (outcome A1). To test whether image
quality alone can produce the under-grading seen on DDR (A2), the development pool was re-scored
out of fold after five fixed degradations of the cached image (JPEG quality 30 at 512 pixels;
downsampling to 256 pixels; Gaussian blur, $\sigma=2$ pixels; gamma 1.8 darkening; contrast
reduced to 60\%), and per-corpus image statistics were computed (A3).

\subsection{Threshold transfer and local recalibration}

A referral threshold on $P(y\geq2)$ was fitted on development predictions for target
sensitivities of 80\%, 85\% (the superiority endpoint of the first autonomous pivotal
trial\cite{abramoff2018pivotal}), 87.2\% and 95.5\%, and applied unchanged externally. To size
local recalibration with the operational objective, $n$ external images ($n=25$--800, 500
draws each) were drawn at random; the referral threshold was set to the largest value giving
at least 85\% sensitivity on those $n$, and sensitivity and specificity were measured on the
remaining images. A draw counts as \emph{harm} when held-out sensitivity falls below 80\%. Positive
and negative predictive values and the number of referable patients missed per 10\,000 screened
were derived at referable prevalences of 5, 10 and 20\%.

\subsection{Statistics}

The primary grading metric is quadratic weighted kappa (QWK)\cite{cohen1968kappa}; the
operationally primary pair is referable sensitivity and specificity, never quoted alone.
Per-grade recall, precision and F1, Cohen's $\kappa$, macro-averaged recall (as a diagnostic,
never an objective), area under the curve, Brier score and expected calibration error are reported
beside them. Intervals are 95\% percentile intervals from 2\,000 bootstrap resamples of groups
(patients where identified, otherwise images); differences between models are paired bootstraps,
and an interval containing zero is reported as indistinguishable, not ranked. Model rankings were
compared with Kendall's $\tau$. A difference between configurations is treated as beyond training
noise only if it exceeds $2\sqrt{2}$ standard deviations of matched development QWK across five
seeds of the reference model. No correction for multiplicity is applied; per-grade claims are
treated as credible only when they replicate across conditions.

\subsection{Software and use of AI tools in the research process}

Every number in the article is generated from archived per-image predictions into a ledger from
which the manuscript is compiled; no number is typed by hand. An AI assistant (Claude, Anthropic)
was used to write analysis code and to run the pre-specified analyses on Kaggle; all code,
analysis plans and outputs are public, and the authors verified the analyses and every cited claim
against its source.
